\documentclass[a4paper]{article}
\input{include/common.tex}

\bibliography{include/references.bib}

\title{%
  Distributed Systems: Lab 1\\%
  Client-Server Database%
}
\author{Ivan Ukhov, Sergiu Rafiliu, Petru Eles, Adrian Horga \\
and Anders Märak Leffler\\
\vspace{0.5em}
}
% \date{December 27, 2016}
\date{}

\begin{document}
\maketitle

\section{Introduction}
\input{include/assets/single-server.tex}

In the previous lab, you built a non-distributed system composed of a number of
standalone machines, and each such machine was responsible for both the server
and client components of your fortune database. You observed that the private
copies of the data maintained by the machines could easily become inconsistent
as there were no attempts of synchronization. The goal of the current lab is to
mitigate this problem by utilizing the client-server model \cite{lecture2} of
distributed systems. Specifically, we shall designate the server role to one
separate machine and turn the rest into pure clients. This scenario is depicted
in \fref{single-server}; compare it with the one from \aref{0}.

\section{Communication Protocol} \slab{protocol}
In order for the server and clients to be able to interact with each other, they
need to agree upon a communication protocol. In this programming project, the
communication protocol has the following specification.

Each message is encoded in the JSON format \cite{json, python-json} and sent as
a single line, which ends with a new-line character. Whenever a client wants to
invoke a function on the server side, the client sends a message to the server
using the following scheme:
\begin{json}
{
    "method": method_name,
    "args": method_arguments
}
\end{json}
The \texttt{method} field contains a string with the name of the function that
the client wants to call, and the \texttt{args} field contains a list of
arguments that that function takes. For example, in order to add a new fortune
into the database of the server, a client might send the following message:
\begin{json}
{
    "method": "write",
    "args": ["Take it easy"]
}
\end{json}
meaning that the method to invoke is called ``write,'' and ``Take it easy''
should be passed as the first and the only argument of that method. As noted
earlier, it is important to send each message as a single line terminated by a
new-line character; make sure you follow this convention when writing into the
socket.

If a method invocation is successful, the server sends an acknowledgment to the
corresponding client using the following scheme:
\begin{json}
{
    "result": method_result
}
\end{json}
where the \texttt{result} field contains the output of the executed function
encoded in the JSON format; the whole message is a single line. Note that the
presence of this field is sufficient to conclude that the call was successful,
and the value of the field can be \code{null} if the corresponding function
returns nothing.

If a method invocation fails due to an error, the server replies with a message
having the following structure:
\begin{json}
{
    "error": {
        "name": error_class_name,
        "args": error_arguments
    }
}
\end{json}

The occurred error is encoded in \texttt{error} in such a way that it
can be instantiated and raised on the client side. In fact, it is one
of the requirements to your implementation.

\subsection{A note on exceptions}

At this point, it is worth recalling that an exception in \python\ is
an instance of some exception class. In order to create such an
instance, one has to know the name of the exception class and the
arguments that the class' constructor requires. With this in mind, the
\texttt{name} field delivers the class name of the occurred exception
while \texttt{args} delivers a list of the parameters to be passed to
the constructor of the class. \emph{In our simplified case, the
  exception class should be a new class inheriting from BaseException,
  and with the right name}.



\section{Your Task}

\subsection{Preparation}
Continue working with the same code base that you have been working
with so far, including all the changes that you have made. The files
relevant to this lab are listed below. You should read and understand
them.
\begin{itemize}

  \item \filename{lab1/client.py} --- the \classname{Client} application (\fix);

  \item \filename{lab1/server.py} --- the \classname{Server} application (\fix);

  \item \filename{lab1/test.sh} --- a shell script that you can use for testing;

  \item \filename{lab1/dbs/fortune.db} --- the same as for \aref{0};

  \item \filename{modules/Common/wrap.sh} --- an auxiliary script needed for
  \filename{test.sh};

  \item \filename{modules/Server/lock/readWriteLock.py} --- the class utilized
  by \classname{Server} for the purpose of synchronization as motivated in
  \sref{communication}. Initially this does nothing. (\fix);

  \item \filename{modules/Server/database.py} --- the same as for \aref{0}
  (\overwrite).

\end{itemize}

Carefully read the source code and understand what it is doing. Pay
attention to the missing parts marked with ``Your code here.''

\subsection{Running the code}
In order to run the code, open two terminal windows and change the
current folder to \filename{lab1}. In one of the terminals, issue the
following command in order to run the \classname{Server} application
(recall that the server is unique in this lab):
\begin{shell}
\$ ./server.py
Listening to: c-204-13.eduroam.liu.se:45725
Press Ctrl-C to stop the server...
\end{shell}
The code randomly chooses a port number to listen to and prints it out along
with the hostname of the server machine. In the example given above, the
hostname is \texttt{c-204-13.eduroam.liu.se}, and the port number is
\texttt{45725}. As you have already discovered by studying the code in
\filename{server.py}, you can choose a port number yourself by specifying the
corresponding argument in the command line. Now, let us create an instance of
the \classname{Client} application by running the following command in the
second terminal:
\begin{shell}
\$ ./client.py -i c-204-13.eduroam.liu.se:45725
Choose one of the following commands:
    r            ::  read a random fortune from the database,
    w <FORTUNE>  ::  write a new fortune into the database,
    h            ::  print this menu,
    q            ::  exit the application.

Command> w Keep it simple.
Command> r
None
Command>
\end{shell}

Note that you have to specify the address that your server is using,
and it can be a different address each time you restart the
server. Regardless of the correctness of the address, the system does
not work at this stage since neither \classname{Server} nor
\classname{Client} has the corresponding communication part
implemented.

In order to ease the testing procedure of this and the future labs, we
have implemented a set of shell scripts that open the needed terminals
and run the corresponding commands for you. These scripts are called
\filename{test.sh} and can be found in the main folder of each
lab. For instance, the test script for the current lab is
\filename{lab1/test.sh}. It is important to note, however, that you
should be able to set up everything yourself without the help of these
scripts. So, have a look at the content of the scripts and make sure
you understand what they do. You can also modify these scripts as you
wish.

\subsection{A note on manual testing}

We use a fairly simple plaintext protocol, where you can easily play
the part of server or client to test out functionality.

Try communicating manually with the server or client part using
``nc``\footnote{See \url{https://wizardzines.com/comics/netcat/} or
  \url{https://linux.die.net/man/1/nc}}. This will allow you to see
how the client tries to communicate with the server, try out sending
malformed data to the client or server to see how they respond, or
respond in a way that should raise an exception at the client side.

As an exercise, start two terminals, one that listens
continuously/repeateadly (with ``-k``) to TCP port 1234. Then try
connecting to it and communicating between the terminals (break with
Ctrl-C).
\begin{shell}
  \$ nc localhost 1234
  Hello
  ^C
\end{shell}

Additionally, try out 
\begin{shell}
  \$ echo '{"method" : "read"}' | nc localhost 1234
\end{shell}


\subsection{Communication} \slab{communication}
Using \python's sockets \cite{python-ipc}, complete the implementation of
\filename{server.py} and \filename{client.py}. Your code should satisfy the
following requirements:
\begin{itemize}

  \item Both \classname{Server} and \classname{Client} should implement and
  follow the communication protocol described in \sref{protocol}.

  \item The server should process each client in a separate thread such that it
  does not stop serving other incoming connections. Your implementation should
  be thread safe as motivated below.

  \item The server should be robust: it should catch and adequately
    respond to exceptions including those due to potential network
    problems, violations of the protocol, invocations of non-existing
    functions, invocations of proper functions but with wrong
    arguments, \etc

  \item Whenever a client receives an error from the server, it should
    instantiate and raise it as if the error has occurred locally. The
    user of the \classname{Client} application should not feel any
    difference between the non-distributed and distributed versions of
    the system. Note that the Python library contains a function for
    dynamic class creation (or, to be more precise, type object
    creation).\cite{python-library}. See above for a discussion of
    what we mean by raising the exception locally.
\end{itemize}

\subsection{Synchronisation}

Since the server can have several threads serving client requests in parallel,
the operations on the database issued by the clients can interfere with each
other and leave the database in a corrupted state. In order to prevent this
behavior, you should protect the critical parts of your code with a read/write
lock provided to you in \filename{readWriteLock.py}. Think where the locking
should take place and how to make it efficient in terms of the waiting time.

Lastly, make sure that, in spite of all possibly occurred exceptions,
the server keeps on working until it is explicitly stopped by the
user. The clients may crash, but the server must keep working.

\subsection{Additional questions}
When demonstrating this lab (as with any lab in the course), you
should be able to explain your code, but also to evaluate it, and
relate it to course concepts. In many cases this will involve
discussing design choices and simplifications that have been made for
you, and may be somewhat speculative. In some cases you might be asked
to answer questions with (at least in retrospect) quite obvious
answers. In this lab, this might include questions such as

\begin{itemize}
  \item What happens if a message is lost, in this design?
  \item In this lab (series) we raise exceptions with the right name,
    that inherit from ``BaseException``, in order to make it seem like
    we are using a local object. Are there any pros and cons of this
    approach?
  \item One cannot focus on everything, and this lab (and course)
    abstracts away many security concerns. It should however be
    mentioned. What security concerns can be raised about our current
    system? Who can access it, and from where?
  \item We use the ReadWriteLock on the server side (in the actual
    database object's read and write methods, in the code
    skeleton). Why don't we provide this on the client side instead? Would
    you notice any difference when doing manual testing with the
    textual user interface?
\end{itemize}

We do not require written answers to this, and you might not need to
provide a solution to all issues raised by the questions.

\section{Conclusion}
In this lab, you have implemented a rather na\"{i}ve distributed system: there
is only one server and each client is required to know the current address of
the server. The former issue makes our distributed system both unreliable (\eg,
your carefully written code will not protect the server against a blackout) and
inefficient (\eg, if there are many clients, the workload of a single server
machine can be substantially high making clients wait for a response for a long
time). The second issue makes it inconvenient for the clients to connect to the
server as the hostname and/or port of the server might easily change, and the
clients need to keep track of such changes. This inconvenience will become
especially bothering when you allow your distributed system to have several
servers; therefore, in the next lab, you will address this problem first.

\printbibliography

\end{document}
